\documentclass[11pt]{article}
\usepackage[margin=1in]{geometry}
\usepackage{lineno}
\usepackage{booktabs}
\usepackage{enumitem}
\usepackage{xcolor}
\usepackage{titling}
\usepackage{titlesec}
\usepackage{tabularx}
\usepackage{multirow}
\usepackage{float}
\usepackage{graphicx}
\usepackage[font=small,labelfont=bf,labelsep=period,
  justification=raggedright,singlelinecheck=false]{caption}
\usepackage[colorlinks=true,linkcolor=black,urlcolor=blue!60!black,citecolor=black]{hyperref}

% Inline references as clickable [N] links (no bibliography section)
\newcommand{\citelink}[2]{\href{#1}{\small[#2]}}
\newcommand{\ok}{\textcolor{green!60!black}{OK}}
\newcommand{\bad}{\textcolor{red!70!black}{MISMATCH}}
\newcommand{\warn}{\textcolor{orange!80!black}{WARN}}

\renewcommand\linenumberfont{\normalfont\scriptsize\sffamily\color[rgb]{0.5,0.5,0.5}}
\leftlinenumbers
\linenumbers

\titleformat{\section}{\large\bfseries}{\thesection.}{0.5em}{}
\titlespacing{\section}{0pt}{2.0em}{0.6em}
\titlespacing{\subsection}{0pt}{1.8em}{0.5em}

\setlength{\droptitle}{-4em}
\pretitle{\begin{center}\large\bfseries}
\posttitle{\end{center}\vspace{-1em}}
\preauthor{\begin{center}\normalsize}
\postauthor{\end{center}\vspace{-1em}}
\predate{\begin{center}\normalsize}
\postdate{\end{center}\vspace{-0.5em}}

\title{Baseline Validation Report}
\author{Diego Toribio\\[0.3em]{\normalsize March 12, 2026}}
\date{\vspace{-2.5em}}

\begin{document}
\maketitle
\thispagestyle{empty}

\section{Purpose}

Before launching batch experiments, we need to confirm that our
implementations produce results in the same ballpark as published
baselines. This report documents:

\begin{enumerate}[leftmargin=2em, itemsep=0.3em]
  \item Current hyperparameters and their source citations.
  \item Discrepancies between our configs and published references.
  \item Validation runs against published numbers.
  \item A change log recording each fix and its measured effect.
\end{enumerate}

This serves as the provenance record for our final hyperparameter
choices and will support the ablation discussion in the thesis.

% =====================================================================
\section{Reference Baselines}
% =====================================================================

All published numbers come from Schwarzer et al.\
\citelink{https://arxiv.org/abs/2007.05929}{1} (SPR paper, Table 4),
which reports mean scores across 10 seeds. Our runs use a single seed
(42), so we expect variance but not order-of-magnitude gaps.

\begin{table}[H]
\centering
\small
\renewcommand{\arraystretch}{1.5}
\setlength{\tabcolsep}{5pt}
\begin{tabular}{l rrrrr}
\toprule
\textbf{Game} & \textbf{Random} & \textbf{Human}
  & \textbf{DER} & \textbf{OTRb} & \textbf{SPR} \\
\midrule
Crazy Climber & 10781 & 35829 & 16185 & 21328 & 42924 \\
Road Runner   & 12    & 7845  & 9600  & 2697  & 14221 \\
Boxing        & 0.1   & 12.1  & 0.2   & 2.5   & 35.8  \\
Kangaroo      & 52    & 3035  & 779   & 605   & 3276  \\
Frostbite     & 65    & 4335  & 867   & 232   & 1822  \\
Up N Down     & 533   & 11693 & 2877  & 2848  & 28139 \\
\bottomrule
\end{tabular}
\caption{Published baselines (10 seeds). DER = Data-Efficient
Rainbow. OTRb = OTRainbow. SPR = SPR with augmentation on Rainbow.
All use Rainbow components.}
\label{tab:published}
\end{table}

% =====================================================================
\section{Current Results vs.\ Published}
% =====================================================================

\subsection{DQN Conditions (seed 42)}

No published vanilla DQN baseline exists for Atari-100K, so these
cannot be validated against external numbers. We include them for
completeness and to establish internal consistency.

\begin{table}[H]
\centering
\small
\renewcommand{\arraystretch}{1.5}
\setlength{\tabcolsep}{6pt}
\begin{tabular}{l rrrr}
\toprule
\textbf{Game} & \textbf{DQN} & \textbf{+ Aug} & \textbf{+ SPR}
  & \textbf{Aug + SPR} \\
\midrule
Crazy Climber & 3540  & 100   & 3400  & 3     \\
Road Runner   & 1440  & 3     & 2003  & 1033  \\
Boxing        & $-3$  & $-16$ & $-17$ & $-40$ \\
Kangaroo      & 27    & 467   & 20    & 320   \\
Frostbite     & 147   & 2     & 157   & 173   \\
Up N Down     & 141   & 1085  & 989   & 491   \\
\bottomrule
\end{tabular}
\caption{Our DQN conditions, single seed. No external baseline
for comparison.}
\end{table}

\subsection{Rainbow Conditions (seed 42)}

\begin{table}[H]
\centering
\small
\renewcommand{\arraystretch}{1.5}
\setlength{\tabcolsep}{6pt}
\begin{tabular}{l rrrl}
\toprule
\textbf{Game} & \textbf{Our Rb} & \textbf{DER} & \textbf{OTRb}
  & \textbf{Gap vs.\ DER} \\
\midrule
Crazy Climber & 557   & 16185 & 21328 & $29\times$ \\
Road Runner   & 850   & 9600  & 2697  & $11\times$ \\
Boxing        & $-12$ & 0.2   & 2.5   & both near-zero \\
Kangaroo      & 40    & 779   & 605   & $20\times$ \\
Frostbite     & 124   & 867   & 232   & $7\times$ \\
Up N Down     & 986   & 2877  & 2848  & $3\times$ \\
\bottomrule
\end{tabular}
\caption{Our Rainbow vs.\ published Rainbow variants. Gaps of
$3$--$29\times$ indicate configuration or implementation issues.}
\label{tab:rainbow-gap}
\end{table}

Our Rainbow underperforms published baselines by an order of magnitude
on most games. Single-seed variance cannot explain a $10$--$29\times$
gap. This points to hyperparameter mismatches, identified in
Section~\ref{sec:audit}.

% =====================================================================
\section{Hyperparameter Audit}
\label{sec:audit}
% =====================================================================

We compare our resolved Rainbow config against the SPR paper's
``controlled Rainbow'' (Table~3, Appendix~A) and the original Rainbow
paper (Hessel et al.\ \citelink{https://arxiv.org/abs/1710.02298}{2},
Appendix~B).

\subsection{Full Comparison}

\begin{table}[H]
\centering
\footnotesize
\renewcommand{\arraystretch}{1.5}
\setlength{\tabcolsep}{3pt}
\begin{tabular}{l l l l l}
\toprule
\textbf{Parameter} & \textbf{Ours} & \textbf{SPR paper}
  & \textbf{Orig.\ Rb} & \textbf{Status} \\
\midrule
Optimizer           & Adam           & Adam              & Adam           & \ok \\
Learning rate       & 6.25e-5        & 1.0e-4            & 6.25e-5        & \warn \\
Adam $\epsilon$     & 1.5e-4         & 1.5e-4            & 1.5e-4         & \ok \\
Train every         & 4 steps        & 1 step            & 4 steps        & \bad \\
Updates / trigger   & 1              & 2                 & 1              & \bad \\
Eff.\ replay ratio  & 0.25           & 2.0               & 0.25           & \bad \\
Multi-step $n$      & 3              & 10                & 3              & \bad \\
Target update       & hard / 8K      & hard / 1 step     & hard / 8K      & \bad \\
Warmup (transitions)& 2000           & 2000              & 80K            & \ok \\
FC hidden dim       & 512            & 256               & 512            & \warn \\
Batch size          & 32             & 32                & 32             & \ok \\
$\gamma$            & 0.99           & 0.99              & 0.99           & \ok \\
PER $\alpha$        & 0.5            & 0.5               & 0.5            & \ok \\
PER $\beta$ anneal  & 0.4 $\to$ 1.0 & 0.4 $\to$ 1.0     & 0.4 $\to$ 1.0 & \ok \\
C51 atoms           & 51             & 51                & 51             & \ok \\
$v_{\min}, v_{\max}$& $-10, 10$      & $-10, 10$         & $-10, 10$      & \ok \\
NoisyNets           & yes            & yes               & yes            & \ok \\
Dueling             & yes            & yes               & yes            & \ok \\
Double DQN          & yes            & yes               & yes            & \ok \\
Loss function       & C51 CE         & C51 CE            & C51 CE         & \ok \\
Grad clip norm      & 10.0           & 10.0              & 10.0           & \ok \\
Reward clipping     & $[-1, 1]$      & $[-1, 1]$         & $[-1, 1]$      & \ok \\
Frame skip          & 4              & 4                 & 4              & \ok \\
Frame stack         & 4              & 4                 & 4              & \ok \\
Eval episodes       & 30             & 100               & --             & \warn \\
\bottomrule
\end{tabular}
\caption{Hyperparameter audit. \ok{} = matches published.
\bad{} = critical mismatch likely causing performance gap.
\warn{} = minor mismatch, lower priority.}
\label{tab:audit}
\end{table}

\subsection{Critical Discrepancies}

Four parameters differ from the SPR paper's data-efficient Rainbow
and are likely responsible for the performance gap:

\medskip
\textbf{1. Training cadence: train every 4 steps vs.\ every 1 step.}
Our config inherits \texttt{train\_every = 4} from the DQN base
(Mnih et al.\ 2015). The SPR paper sets ``Replay period every: 1
step'' (Table~3), meaning a training trigger fires on every
environment step. Combined with the replay ratio difference below,
this affects both the total number of gradient updates and how
frequently fresh experience enters the training loop.

Source: Schwarzer et al.\ (2021), Table~3: ``Replay period every =
1 step''.

\medskip
\textbf{2. Replay ratio: 0.25 vs.\ 2.0 (8$\times$ fewer gradient
steps).} With \texttt{train\_every = 4} and 1 update per trigger, we
get 0.25 gradient updates per environment step. The SPR paper
performs 2 updates per 1 step. Over 100K steps with warmup = 2000,
SPR accumulates $\sim$196K gradient updates; we accumulate
$\sim$25K---nearly an order of magnitude fewer.

Together, items 1 and 2 define the training intensity. To match the
SPR paper exactly: \texttt{train\_every = 1},
\texttt{updates\_per\_step = 2}, giving 2 gradient updates per
environment step.

Source: Schwarzer et al.\ (2021), Table~3: ``Updates per step = 2''.

\medskip
\textbf{3. Multi-step $n$: 3 vs.\ 10.} We use the original Rainbow
default ($n=3$). The SPR paper uses $n=10$, propagating reward signal
10 steps back through trajectories. In the 100K regime where episodes
are short and data scarce, longer n-step returns provide
substantially better learning signal.

Source: Schwarzer et al.\ (2021), Table~3: ``multi-step return n = 10''.

\medskip
\textbf{4. Target network: hard copy every 8K steps vs.\ every 1
step.} With our replay ratio of 0.25, we perform $\sim$25K gradient
steps total. A hard update every 8K steps means roughly 3 target
updates during the entire run. The SPR paper updates every step,
keeping targets fresh.

Source: Schwarzer et al.\ (2021), Table~3: ``target network update
period = 1''.

\subsection{Minor Discrepancies}

These are lower priority but should be aligned after fixing the
critical items:

\begin{itemize}[leftmargin=2em, itemsep=0.3em]
  \item \textbf{Learning rate}: 6.25e-5 vs.\ 1.0e-4 (SPR paper uses
    higher LR with higher replay ratio).
  \item \textbf{FC hidden dim}: 512 vs.\ 256 (larger network may
    overfit at high replay ratio; the SPR paper's 256 was tuned for
    replay ratio 2).
  \item \textbf{Eval episodes}: 30 vs.\ 100. The Atari-100K
    benchmark (Kaiser et al.\ 2020) specifies 100 evaluation
    trajectories; SPR follows this. Our Rainbow config uses Noisy
    Nets for exploration (no eval $\epsilon$), but we currently
    evaluate only 30 episodes.
\end{itemize}

\subsection{Root Cause}

Our Rainbow config inherits directly from the original Rainbow paper
(Hessel et al.\ 2018), which was designed for 200M-frame training
runs. The hyperparameters that differ (replay ratio 0.25, $n$-step 3,
target update every 8K) are all sensible defaults for long training
but catastrophically underpowered for the 100K-step regime.

This diagnosis is confirmed by Kielak
\citelink{https://arxiv.org/abs/2003.10181}{3} (OTRainbow), whose
``SRainbow'' column (standard Rainbow with Dopamine defaults at 100K
steps) reports scores in the same range as ours---e.g., Crazy Climber
2254 (SRainbow) vs.\ 557 (ours), both far below tuned DER at 16185.
The gap is purely hyperparameter tuning, not an implementation bug.

The SPR paper adapted these parameters specifically for data
efficiency: higher replay ratio extracts more signal from limited
data, longer n-step returns propagate reward faster, and frequent
target updates prevent staleness when the network changes rapidly.

% =====================================================================
\section{DQN Hyperparameter Record}
\label{sec:dqn-params}
% =====================================================================

Our DQN conditions cannot be compared to published baselines (no
vanilla DQN exists for Atari-100K), but we record the parameters
and their sources for completeness.

\begin{table}[H]
\centering
\small
\renewcommand{\arraystretch}{1.5}
\setlength{\tabcolsep}{5pt}
\begin{tabular}{l l l}
\toprule
\textbf{Parameter} & \textbf{Value} & \textbf{Source} \\
\midrule
Architecture   & Nature CNN (32-64-64-512) & Mnih et al.\ (2015) \\
Optimizer      & RMSProp, lr=2.5e-4        & Mnih et al.\ (2015) \\
Batch size     & 32                        & Mnih et al.\ (2015) \\
Replay capacity& 100K                      & Matched to 100K budget \\
Warmup         & 2K transitions            & SPR paper default \\
Target update  & Hard copy every 2K steps  & Scaled from 10K/10M \\
$\epsilon$ decay& 1.0 $\to$ 0.1 / 200K frames & First half of training \\
$\gamma$       & 0.99                      & Mnih et al.\ (2015) \\
Grad clip      & L2 norm $\leq$ 10         & Standard practice \\
Train every    & 4 steps                   & Mnih et al.\ (2015) \\
Eval episodes  & 30                        & SPR paper protocol \\
Eval $\epsilon$& 0.05                      & Mnih et al.\ (2015) \\
\bottomrule
\end{tabular}
\caption{DQN hyperparameters with source citations. These were not
modified from initial implementation.}
\label{tab:dqn-params}
\end{table}

\subsection{SPR-Specific Parameters (DQN + SPR)}

\begin{table}[H]
\centering
\small
\renewcommand{\arraystretch}{1.5}
\setlength{\tabcolsep}{5pt}
\begin{tabular}{l l l}
\toprule
\textbf{Parameter} & \textbf{Value} & \textbf{Source} \\
\midrule
Prediction steps $K$ & 5            & SPR paper, Table 3 \\
Loss weight $\lambda$& 2.0          & SPR paper, Table 3 \\
Projection dim       & 512          & Matches FC width \\
Transition model     & 2 conv, 64ch & SPR paper, Section 2.3 \\
EMA momentum $\tau$  & 0.99         & SPR paper (no aug) \\
Dropout              & 0.5          & SPR paper (no aug) \\
\bottomrule
\end{tabular}
\caption{SPR parameters. With augmentation: $\tau = 0.0$ (direct
copy), dropout = 0.0.}
\label{tab:spr-params}
\end{table}

% =====================================================================
\section{Fix Plan}
\label{sec:fix-plan}
% =====================================================================

We fix discrepancies in priority order and re-run on 2 validation
games (Crazy Climber and Up N Down---chosen for large
published-vs-ours gap and dense rewards) to measure the delta.

\medskip
\begin{enumerate}[leftmargin=2em, itemsep=0.5em]
  \item \textbf{Training cadence + replay ratio + multi-step +
    target update}: Set \texttt{train\_every = 1},
    \texttt{updates\_per\_step = 2} (replay ratio 2.0, matching SPR
    Table~3 exactly), \texttt{multi\_step.n = 10},
    \texttt{target\_update = 1}, \texttt{warmup\_steps = 2000},
    \texttt{min\_size = 2000}. This brings total gradient updates
    from $\sim$25K to $\sim$196K and starts training 18K steps
    earlier.
  \item \textbf{Learning rate}: Change from 6.25e-5 to 1.0e-4.
  \item \textbf{FC hidden}: Reduce from 512 to 256 in Rainbow model
    (requires RainbowDQN architecture change).
\end{enumerate}

Step 1 contains all critical fixes applied together in the first
validation run (these parameters are interdependent---replay ratio
and learning rate were co-tuned in the SPR paper). Steps 2--3 will
be applied individually to isolate their effects.

% =====================================================================
\section{Change Log}
\label{sec:changelog}
% =====================================================================

Each entry records the change, the validation game(s), and the
before/after scores. This section will be filled as fixes are
implemented and tested.

\subsection{Baseline (pre-fix)}

\begin{table}[H]
\centering
\small
\renewcommand{\arraystretch}{1.5}
\setlength{\tabcolsep}{6pt}
\begin{tabular}{l rr}
\toprule
\textbf{Game} & \textbf{Our Rb} & \textbf{DER (published)} \\
\midrule
Crazy Climber & 557   & 16185 \\
Up N Down     & 986   & 2877  \\
\bottomrule
\end{tabular}
\caption{Pre-fix Rainbow scores on validation games (seed 42).}
\end{table}

\subsection{Fix 1: Training cadence + replay ratio + multi-step +
target update + warmup}

\noindent\textit{Pending.}

\noindent Changes:
\begin{itemize}[leftmargin=2em, itemsep=0.1em]
  \item \texttt{train\_every = 1}, \texttt{updates\_per\_step = 2}
    (replay ratio 2.0)
  \item \texttt{multi\_step.n = 10}, \texttt{target\_update = 1}
  \item \texttt{warmup\_steps = 2000}, \texttt{min\_size = 2000}
\end{itemize}

\noindent\textit{Note:} Previous configs used
\texttt{train\_every\,=\,4} with
\texttt{updates\_per\_step\,=\,8}, which gives the same effective
ratio but different training dynamics (bursty 8-update blocks every
4~steps vs.\ smooth 2-update blocks every step). The new settings
match SPR Table~3 exactly.

\subsection{Fix 2: Learning rate}

\noindent\textit{Pending.}\\
\noindent Change: \texttt{lr = 1.0e-4}.

\subsection{Fix 3: FC hidden dimension}

\noindent\textit{Pending.}\\
\noindent Change: \texttt{fc\_hidden = 256} for Rainbow model only.

% =====================================================================
\section{Acceptance Criteria}
% =====================================================================

Rainbow is ``validated'' when our single-seed scores on the 2
validation games fall within $2\times$ of published DER numbers (a
conservative bar given single-seed variance). Specifically:

\begin{itemize}[leftmargin=2em, itemsep=0.3em]
  \item Crazy Climber: $> 8000$ (vs.\ DER 16185)
  \item Up N Down: $> 1400$ (vs.\ DER 2877)
\end{itemize}

Once validated, we lock the Rainbow config and proceed to batch runs.

\end{document}
